\section{گام دوم: طراحی معماری شبکه‌های عصبی}

در این بخش، سه معماری مختلف برای انجام عملیات طبقه‌بندی و ایجاد یک بستر مقایسه‌ی عادلانه \lr{(Fair Comparison)} توسعه داده شد. هدف، مقایسه‌ی عملکرد یادگیری انتقالی \lr{(Transfer Learning)} با مدلی است که از پایه آموزش می‌بیند.

\subsection{طراحی شبکه‌ی کانولوشنی اختصاصی \lr{(Custom CNN)}}
شبکه‌ی پایه از صفر \lr{(Scratch)} طراحی شد تا به‌عنوان نقطه‌ی شروع \lr{(Baseline)} عمل کند. این معماری شامل سه بلوک کانولوشنی است. در هر بلوک از لایه‌های کانولوشن دو‌بعدی، نرمال‌سازی دسته‌ای \lr{(Batch Normalization)} برای تسریع همگرایی، تابع فعال‌ساز \lr{ReLU} و در نهایت تجمیع حداکثری \lr{(Max Pooling)} استفاده شده است.

\begin{tipbox}[کاهش پارامترها و جلوگیری از بیش‌برازش]
برای اتصال لایه‌های استخراج ویژگی به لایه‌ی تمام‌متصل \lr{(Fully Connected)} در انتهای شبکه، به‌جای استفاده مستقیم از لایه‌ی \lr{Flatten} که باعث افزایش شدید تعداد پارامترها می‌شود، از \lr{Adaptive Average Pooling} استفاده شد تا ابعاد فضایی به $1 \times 1$ کاهش یابد. این کار خطر بیش‌برازش \lr{(Overfitting)} را به شدت کاهش می‌دهد. همچنین از لایه‌ی \lr{Dropout} با نرخ ۰.۵ پیش از طبقه‌بند نهایی استفاده گردید.
\end{tipbox}

همچنین با توجه به استفاده از تابع زیان \lr{CrossEntropyLoss} در مراحل بعد، از قرار دادن لایه‌ی \lr{Softmax} در خروجی شبکه خودداری شد، چرا که این تابع زیان، مقادیر خام \lr{(Logits)} را می‌پذیرد.

\subsection{پیاده‌سازی مدل یادگیری انتقالی \lr{(Transfer Learning)}}
برای مدل از پیش آموزش‌دیده، معماری \lr{ResNet18} از کتابخانه‌ی \lr{Torchvision} انتخاب گردید. دلایل این انتخاب عبارتند از:
\begin{itemize}
    \item \textbf{پیچیدگی محاسباتی متعادل:} این شبکه با وجود عمق مناسب، تعداد پارامترهای معقولی دارد (حدود ۱۱ میلیون پارامتر) که آموزش آن را روی سخت‌افزارهای محدود امکان‌پذیر می‌کند.
    \item \textbf{سرعت استنتاج \lr{(Inference Time)}:} استفاده از اتصالات میان‌بر \lr{(Skip Connections)} در معماری \lr{ResNet} باعث بهینه‌سازی گرادیان‌ها و سرعت آموزش می‌شود.
\end{itemize}

در این استراتژی، تمام لایه‌های پایه‌ی شبکه \lr{(Backbone)} فریز شدند (\lr{requires\_grad = False}) و تنها لایه‌ی \lr{Fully Connected} نهایی با یک لایه‌ی جدید (متناسب با ۷ کلاس احساسات چهره) جایگزین شد.

\subsection{مدل \lr{ResNet18} با وزن‌دهی تصادفی \lr{(Scratch)}}
برای ایجاد شرایط مقایسه‌ی کاملاً عادلانه، همان معماری \lr{ResNet18} این بار بدون وزن‌های اولیه‌ی \lr{ImageNet} بارگذاری شد. تمامی لایه‌های این مدل قابلیت آموزش دارند. این کار به ما نشان خواهد داد که بهبود عملکرد در مدل قبلی، دقیقاً ناشی از وزن‌های از پیش‌ آموزش‌دیده است یا صرفاً پیچیدگی خود معماری \lr{ResNet18}.

\subsection{تحلیل تعداد پارامترهای مدل‌ها}
تعداد پارامترهای هر سه مدل محاسبه و در ادامه مقایسه شده‌اند. با توجه به فریز شدن لایه‌ها در یادگیری انتقالی، تفاوت چشمگیری در پارامترهای قابل آموزش وجود دارد:

\begin{itemize}
    \item \textbf{مدل \lr{Custom CNN}:} \\ 
    کل پارامترها: $102,407$ | قابل آموزش: $102,407$ | فریز شده: $0$
    
    \item \textbf{مدل \lr{Pretrained ResNet18 (Transfer Learning)}:} \\ 
    کل پارامترها: $11,180,103$ | قابل آموزش: $3,591$ | فریز شده: $11,176,512$
    
    \item \textbf{مدل \lr{Scratch ResNet18}:} \\ 
    کل پارامترها: $11,180,103$ | قابل آموزش: $11,180,103$ | فریز شده: $0$
\end{itemize}

همان‌طور که مشاهده می‌شود، در روش \lr{Transfer Learning}، با فریز کردن استخراج‌گر ویژگی، تعداد پارامترهای قابل آموزش به تنها ۳,۵۹۱ پارامتر (مربوط به طبقه‌بند نهایی) کاهش یافته است. این کاهش چشمگیر در تعداد پارامترهای درگیر در محاسبه‌ی گرادیان، باعث کاهش محسوس زمان آموزش در هر دوره \lr{(Epoch)} و همچنین کاهش نیاز به حافظه‌ی کارت گرافیک \lr{(VRAM)} خواهد شد.